% Generated draft: requires mathematical and source review.
\documentclass[11pt]{article}
\usepackage[margin=25mm]{geometry}
\usepackage{fontspec}
\setmainfont{Latin Modern Roman}
\usepackage{amsmath,amssymb,mathtools,unicode-math}
\setmathfont{Latin Modern Math}
\usepackage{enumitem,needspace,xurl,hyperref}
\setlength{\parindent}{0pt}
\setlength{\parskip}{5pt}
\setlength{\emergencystretch}{3em}
\newcommand{\N}{\mathbb{N}}
\newcommand{\Z}{\mathbb{Z}}
\newcommand{\Q}{\mathbb{Q}}
\newcommand{\R}{\mathbb{R}}
\newcommand{\C}{\mathbb{C}}
\newcommand{\F}{\mathbb{F}}
\newcommand{\A}{\mathbb{A}}
\newcommand{\PP}{\mathbb{P}}
\newcommand{\E}{\mathbb{E}}
\newcommand{\eps}{\varepsilon}
\newcommand{\dd}{\,\mathrm{d}}
\begin{document}
\section*{UnsolvedMath ID: 30004446}
\textbf{Source problem code:} catalogue entry
\subsection{Definitions and mathematical statement}
Put \(I=[0,1]\). We use the conventional interpretation of \(\mathrm{PL}_o(I)\) as the group, under composition, of orientation-preserving piecewise-linear homeomorphisms of \(I\) having finitely many linear pieces. Such homeomorphisms fix both endpoints. Let Thompson's group \(F\leq \mathrm{PL}_o(I)\) consist of those maps whose breakpoints are dyadic rationals and whose slopes belong to \(\{2^n:n\in\Z\}\). These are standard interpretations of the group names in the supplied record; the record itself does not specify these conventions.

For groups \(G,H\), write
\[
G\preccurlyeq H
\quad\Longleftrightarrow\quad
\text{there exists an injective group homomorphism }G\longrightarrow H.
\]
This is abstract group embeddability, without any stipulated compatibility with the given interval actions. Let \(\mathcal C_F\) and \(\mathcal C_{\mathrm{PL}}\) be the collections of isomorphism classes of finitely generated subgroups of \(F\) and \(\mathrm{PL}_o(I)\), respectively. Equip both collections with the quasiorder \(\preccurlyeq\). A subcollection \(\mathcal D\subseteq\mathcal C\) is an initial segment if
\[
H\in\mathcal D,\quad G\in\mathcal C,\quad G\preccurlyeq H
\quad\Longrightarrow\quad G\in\mathcal D.
\]
Call \(G\in\mathcal C\) universally comparable within \(\mathcal C\) if
\[
\forall H\in\mathcal C,\qquad
G\preccurlyeq H\ \text{or}\ H\preccurlyeq G.
\]
Let \(\mathrm{EA}\) denote the class of elementary amenable groups: the smallest class containing all finite groups and all abelian groups and closed under taking subgroups, quotients, group extensions, and directed unions.

\begin{enumerate}
\item For each of \((\mathcal C_F,\preccurlyeq)\) and \((\mathcal C_{\mathrm{PL}},\preccurlyeq)\), identify a substantial classifiable initial segment and determine where classification becomes intractable. The supplied record gives no mathematical definitions of ``substantial'', ``classifiable'', or ``intractable'', and specifies no classification invariant or complexity framework. Consequently this part remains a research programme rather than a determinate proposition.
\item Determine the universally comparable members of the relevant collection. In particular, is
\[
\forall H\in\mathcal C_{\mathrm{PL}},\qquad
F\preccurlyeq H\ \text{or}\ H\preccurlyeq F?
\]
The source does not explicitly identify the ambient collection for universal comparability. Both collections above are therefore retained as possible scopes. Within \(\mathcal C_F\), comparability with \(F\) follows directly from the definition of \(\mathcal C_F\); the displayed question records the broader, nontrivial interpretation without claiming that the source resolves this scope ambiguity.
\item Is the following implication true for every subgroup \(G\leq\mathrm{PL}_o(I)\)?
\[
F\not\preccurlyeq G\quad\Longrightarrow\quad G\in\mathrm{EA}.
\]
Here ``not containing \(F\)'' is interpreted as containing no subgroup isomorphic to \(F\). This clause of the supplied statement says ``a subgroup'', so the displayed quantifier includes subgroups that are not finitely generated. The opening restriction to finitely generated subgroups leaves it uncertain whether that restriction was also intended for this clause.
\item Is \((\mathcal C_F,\preccurlyeq)\) well quasiordered? Explicitly, does every sequence \((G_n)_{n\in\N}\) of finitely generated subgroups of \(F\) satisfy
\[
\exists i,j\in\N\qquad i<j\quad\text{and}\quad G_i\preccurlyeq G_j?
\]
\end{enumerate}

These questions retain the supplied open-problem framing. Their current status and the unresolved scope conventions require human mathematical and source review.
\subsection{Short English statement}
The record asks how finitely generated subgroups of Thompson's group \(F\) and of the orientation-preserving piecewise-linear interval group compare under injective homomorphisms. It also asks whether avoiding a copy of \(F\) forces elementary amenability, and whether every infinite sequence of finitely generated subgroups of \(F\) has an earlier group embedding into a later one. The classification criteria and some quantifier scopes remain unspecified. The supplied open status is dated curation, not an independently verified conclusion.
\subsection{Sources}
\begin{itemize}
\item Supplied local UnsolvedMath registry record, permanent ID \(30004446\), title: \emph{Embedding Classification of Piecewise-Linear Interval Groups}; collection: Oberwolfach Reports Open Problems; exact source problem code: \texttt{catalogue entry}. No individual proposer is supplied. The recorded proposed year is 2020.
\item Supplied source citation: \emph{Set Theory (2021)}, Oberwolfach Reports, pp.\ 797--855, DOI \texttt{10.4171/owr/2020/14}. The source URL and the literature-source entry labelled ``Oberwolfach Report 14/2020'' both give \url{https://doi.org/10.4171/owr/2020/14}. The supplied title year 2021 and report designation 14/2020 are preserved without resolving their relationship. No problem-specific page or speaker attribution is supplied.
\item Supplied literature-source entry: ``Hyde--Moore (2023)'', \url{https://doi.org/10.4171/GGD/708}. The record supplies no article title or precise theorem reference.
\item Literature assessment checked at \texttt{2026-08-22}: this exact date records supplied dated curation, not an independently verified literature search. The supplied status is ``open''. The supplied assessment states that the general \(F\)-obstruction converse is unknown and that later work gives strong sufficient obstructions without providing the requested classification. The background explicitly calls for independent verification of time-sensitive status and cited claims.
\item The local statement-status label is \texttt{corrected\_verified}; its accompanying qualification says that only formula delimiters, notation, grammar, and typography were repaired while mathematical content and scope were preserved. This label does not establish independent mathematical verification of this draft or of the literature assessment.
\end{itemize}
\end{document}
